\begin{abstract}
Hallucination in Large Language Models (LLMs), where models generate responses that appear convincing but are factually incorrect, remains a major challenge affecting the reliability of AI systems. Measuring hallucinations is an important first step toward addressing this issue; however, most existing evaluation methods rely on benchmark datasets containing human-written gold-standard answers. Creating such datasets is expensive and can introduce human bias or errors. To address this limitation, this study introduces \textbf{Factualness Evaluations via Weighting LLMs (FEWL)}, a new metric designed to measure hallucinations even when gold-standard answers are not available. The proposed approach uses responses generated by multiple existing LLMs as references and assigns them different weights based on their estimated reliability. Experimental results demonstrate that FEWL provides more accurate hallucination measurements than naive reference-based approaches. Furthermore, the metric can also be used to reduce hallucinations through techniques such as in-context learning and supervised fine-tuning.
\end{abstract}